% GENERATED FILE. Edit canonical lessons, then run npm run generate:books.
% canonical-source-hash: fnv1a64:9a2cf08baf6c9c8e
% canonical-lessons: TE-C538-svaram, TE-C538-gatram, TE-C538-vayincu, TE-C538-mrdangam, TE-C538-sannayi

\chapter{Music}
\label{ch:te-music}
\hlchaptermodality{538}

\begin{chapteropening}
\textbf{By the end of this chapter:} \emph{I can talk about notes, voices and the instruments people play.}

Complete the last lesson of chapter 538: I can talk about notes, voices and the instruments people play.
\end{chapteropening}

\section[\texorpdfstring{\te{స్వరం} (svaraṁ)}{స్వరం (svaraṁ)}]{\te{స్వరం} (svaraṁ) — a musical note; a voice}

\label{lesson:TE-C538-svaram}



\begin{quote}
Before the new one: say the Telugu for kabaddi, then the Telugu for to hop.
\end{quote}

\subsection*{You'll want to know: \te{స్వరం}}
\textbf{\te{స్వరం}} — \emph{svaraṁ} — \textquotedblleft{}a musical note; a voice\textquotedblright{}.

\begin{quote}
\te{ఆమె} \te{స్వరం} \te{చాలా} \te{మధురంగా} \te{ఉంటుంది}.
\end{quote}

\begin{quote}
\emph{āme svaraṁ cālā madhuraṅgā uṇṭundi.}
\end{quote}

\begin{quote}
Her voice is very sweet.
\end{quote}

Plural: \te{స్వరాలు}, \emph{svarālu}.

Plural \te{స్వరాలు}: the seven notes of Indian music.

\begin{cousinweb}
\textbf{\te{స్వరం}} (\emph{svaraṁ}) is borrowed from Sanskrit \emph{svara}, \textquotedblleft{}sound; a musical note\textquotedblright{} (CDIAL 13908).
\end{cousinweb}

\begin{grammarlens}[title={What you've built}]
One more word for describing what you lived and felt.
\end{grammarlens}

\subsection*{Your turn}
\begin{itemize}
\raggedright
  \item \emph{Say it:} \emph{svaraṁ}
  \item \emph{Say it:} \emph{svaraṁ}, once more
  \item \emph{Say it:} \emph{āme svaraṁ cālā madhuraṅgā uṇṭundi.} — Her voice is very sweet.
  \item \emph{Recall:} say the Telugu for kabaddi, then the Telugu for to hop, then say \emph{svaraṁ} again
\end{itemize}

\begin{tcolorbox}[breakable,colback=teal!4,colframe=teal!35!black,title={Before you move on}]
What does \te{స్వరం} mean? (\textquotedblleft{}A musical note; a voice\textquotedblright{}.) Say it once more.
\end{tcolorbox}
\section[\texorpdfstring{\te{గాత్రం} (gātraṁ)}{గాత్రం (gātraṁ)}]{\te{గాత్రం} (gātraṁ) — a singing voice}

\label{lesson:TE-C538-gatram}



\begin{quote}
Before the new one: say the Telugu for to hop, then the Telugu for a musical note.
\end{quote}

\subsection*{You'll want to know: \te{గాత్రం}}
\textbf{\te{గాత్రం}} — \emph{gātraṁ} — \textquotedblleft{}a singing voice\textquotedblright{}.

\begin{quote}
\te{ఆమె} \te{గాత్రం} \te{విని} \te{అందరూ} \te{ఆశ్చర్యపోయారు}.
\end{quote}

\begin{quote}
\emph{āme gātraṁ vini andarū āścaryapōyāru.}
\end{quote}

\begin{quote}
Everyone was amazed to hear her singing voice.
\end{quote}

In music, singing as against instruments.

\begin{cousinweb}
\textbf{\te{గాత్రం}} (\emph{gātraṁ}) is borrowed from Sanskrit \emph{gātra}, \textquotedblleft{}a limb; the body\textquotedblright{} (CDIAL 4124); in Telugu music it is the voice, the body's own instrument.
\end{cousinweb}

\begin{grammarlens}[title={What you've built}]
One more word for describing what you lived and felt.
\end{grammarlens}

\subsection*{Your turn}
\begin{itemize}
\raggedright
  \item \emph{Say it:} \emph{gātraṁ}
  \item \emph{Say it:} \emph{gātraṁ}, once more
  \item \emph{Say it:} \emph{āme gātraṁ vini andarū āścaryapōyāru.} — Everyone was amazed to hear her singing voice.
  \item \emph{Recall:} say the Telugu for to hop, then the Telugu for a musical note, then say \emph{gātraṁ} again
\end{itemize}

\begin{tcolorbox}[breakable,colback=teal!4,colframe=teal!35!black,title={Before you move on}]
What does \te{గాత్రం} mean? (\textquotedblleft{}A singing voice\textquotedblright{}.) Say it once more.
\end{tcolorbox}
\section[\texorpdfstring{\te{వాయించు} (vāyiñcu)}{వాయించు (vāyiñcu)}]{\te{వాయించు} (vāyiñcu) — to play an instrument}

\label{lesson:TE-C538-vayincu}



\begin{quote}
Before the new one: say the Telugu for a musical note, then the Telugu for a singing voice.
\end{quote}

\subsection*{You'll want to know: \te{వాయించు}}
\textbf{\te{వాయించు}} — \emph{vāyiñcu} — \textquotedblleft{}to play an instrument\textquotedblright{}.

\begin{quote}
\te{మా} \te{అన్నయ్య} \te{గిటారు} \te{బాగా} \te{వాయిస్తాడు}.
\end{quote}

\begin{quote}
\emph{mā annayya giṭāru bāgā vāyistāḍu.}
\end{quote}

\begin{quote}
My elder brother plays the guitar well.
\end{quote}

Past (I …): \te{వాయించాను}, \emph{vāyiñcānu}.

For instruments; \te{ఆడు} is for games.

\begin{cousinweb}
\textbf{\te{వాయించు}} (\emph{vāyiñcu}) is usually traced to Sanskrit \emph{vādayati}, \textquotedblleft{}causes to sound\textquotedblright{} (CDIAL 11510), with the Telugu verb ending -\te{ఇంచు}.
\end{cousinweb}

\begin{grammarlens}[title={What you've built}]
One more word for describing what you lived and felt.
\end{grammarlens}

\subsection*{Your turn}
\begin{itemize}
\raggedright
  \item \emph{Say it:} \emph{vāyiñcu}
  \item \emph{Say it:} \emph{vāyiñcu}, once more
  \item \emph{Say it:} \emph{mā annayya giṭāru bāgā vāyistāḍu.} — My elder brother plays the guitar well.
  \item \emph{Recall:} say the Telugu for a musical note, then the Telugu for a singing voice, then say \emph{vāyiñcu} again
\end{itemize}

\begin{tcolorbox}[breakable,colback=teal!4,colframe=teal!35!black,title={Before you move on}]
What does \te{వాయించు} mean? (\textquotedblleft{}To play an instrument\textquotedblright{}.) Say it once more.
\end{tcolorbox}
\section[\texorpdfstring{\te{మృదంగం} (mṛdaṅgaṁ)}{మృదంగం (mṛdaṅgaṁ)}]{\te{మృదంగం} (mṛdaṅgaṁ) — a mridangam drum}

\label{lesson:TE-C538-mrdangam}



\begin{quote}
Before the new one: say the Telugu for a singing voice, then the Telugu for to play an instrument.
\end{quote}

\subsection*{You'll want to know: \te{మృదంగం}}
\textbf{\te{మృదంగం}} — \emph{mṛdaṅgaṁ} — \textquotedblleft{}a mridangam drum\textquotedblright{}.

\begin{quote}
\te{పాటకు} \te{తోడుగా} \te{ఆయన} \te{మృదంగం} \te{వాయించారు}.
\end{quote}

\begin{quote}
\emph{pāṭaku tōḍugā āyana mṛdaṅgaṁ vāyiñcāru.}
\end{quote}

\begin{quote}
He played the mridangam to accompany the song.
\end{quote}

A two-headed barrel drum of South Indian music.

\begin{cousinweb}
\textbf{\te{మృదంగం}} (\emph{mṛdaṅgaṁ}) is borrowed from Sanskrit \emph{mṛdaṁga}, \textquotedblleft{}a kind of drum\textquotedblright{} (Monier-Williams).
\end{cousinweb}

\begin{grammarlens}[title={What you've built}]
One more word for describing what you lived and felt.
\end{grammarlens}

\subsection*{Your turn}
\begin{itemize}
\raggedright
  \item \emph{Say it:} \emph{mṛdaṅgaṁ}
  \item \emph{Say it:} \emph{mṛdaṅgaṁ}, once more
  \item \emph{Say it:} \emph{pāṭaku tōḍugā āyana mṛdaṅgaṁ vāyiñcāru.} — He played the mridangam to accompany the song.
  \item \emph{Recall:} say the Telugu for a singing voice, then the Telugu for to play an instrument, then say \emph{mṛdaṅgaṁ} again
\end{itemize}

\begin{tcolorbox}[breakable,colback=teal!4,colframe=teal!35!black,title={Before you move on}]
What does \te{మృదంగం} mean? (\textquotedblleft{}A mridangam drum\textquotedblright{}.) Say it once more.
\end{tcolorbox}
\section[\texorpdfstring{\te{సన్నాయి} (sannāyi)}{సన్నాయి (sannāyi)}]{\te{సన్నాయి} (sannāyi) — a reed pipe; a shehnai}

\label{lesson:TE-C538-sannayi}



\begin{quote}
Before the new one: say the Telugu for to play an instrument, then the Telugu for a mridangam drum.
\end{quote}

\subsection*{You'll want to know: \te{సన్నాయి}}
\textbf{\te{సన్నాయి}} — \emph{sannāyi} — \textquotedblleft{}a reed pipe; a shehnai\textquotedblright{}.

\begin{quote}
\te{పెళ్ళిలో} \te{సన్నాయి} \te{మోగింది}.
\end{quote}

\begin{quote}
\emph{peḷḷilō sannāyi mōgindi.}
\end{quote}

\begin{quote}
The reed pipe played at the wedding.
\end{quote}

The festive reed instrument of weddings.

\begin{cousinweb}
\textbf{\te{సన్నాయి}} (\emph{sannāyi}) is the Telugu name of the festive reed pipe; its origin is not settled.
\end{cousinweb}

\begin{grammarlens}[title={What you've built}]
One more word for describing what you lived and felt.
\end{grammarlens}

\subsection*{Your turn}
\begin{itemize}
\raggedright
  \item \emph{Say it:} \emph{sannāyi}
  \item \emph{Say it:} \emph{sannāyi}, once more
  \item \emph{Say it:} \emph{peḷḷilō sannāyi mōgindi.} — The reed pipe played at the wedding.
  \item \emph{Recall:} say the Telugu for to play an instrument, then the Telugu for a mridangam drum, then say \emph{sannāyi} again
\end{itemize}

\begin{tcolorbox}[breakable,colback=teal!4,colframe=teal!35!black,title={Before you move on}]
What does \te{సన్నాయి} mean? (\textquotedblleft{}A reed pipe; a shehnai\textquotedblright{}.) Say it once more.
\end{tcolorbox}
